\section{Pairs of vertices: the grid picture}\label{sec:pairs}

\begin{idea}[What this section does]
We introduce the language of the rest of the proof: subsets of the grid
$X\times X$, the two coordinate images $\Ff,\Fs$, conflicts and compatibility,
and the quantity $\val(P,Q)$ that will be maximized. We then explain
\emph{heuristically} why maximizing $\val$ leads to a contradiction. The
explanation is made rigorous in \cref{sec:contradiction}, and the one
missing tool, the pruning lemma, is proved in \cref{sec:pruning}.
\end{idea}

\subsection{Coordinate images}

In this section and the next, $X$ is a finite set with an arbitrary binary
relation $\to$. Loops $x\to x$ and pairs of opposite arrows are
allowed; this generality is needed for the pruning lemma. In
\cref{sec:contradiction} the relation will be the arc relation of the graph
$G$ from \cref{prop:reduction}. We write $F(U)=\{y: x\to y\text{ for some }x\in U\}$
as before.

\emph{Drawing convention.} A pair $(i,j)\in X\times X$ is drawn in
\textbf{row $i$ and column $j$} of a square grid. So the first coordinate is
vertical and the second is horizontal.

\begin{definition}[Coordinate images, columns, rows]\label{def:coord-images}
For $S\subseteq X\times X$ let
\begin{align*}
 \Ff S&=\{(i,j):\ (p,j)\in S\text{ and }p\to i\text{ for some }p\}
 &&\text{(move the first coordinate),}\\
 \Fs S&=\{(i,j):\ (i,s)\in S\text{ and }s\to j\text{ for some }s\}
 &&\text{(move the second coordinate),}
\end{align*}
and $\Ff^2=\Ff\circ\Ff$, $\Fs^2=\Fs\circ\Fs$. The \emph{column} $j$ of $P$ and the
\emph{row} $i$ of $Q$ are
\[
 P_j=\{p:(p,j)\in P\},\qquad Q_i=\{s:(i,s)\in Q\}.
\]
\end{definition}

So $\Ff$ moves points \emph{vertically} inside their column and $\Fs$ moves
points \emph{horizontally} inside their row (\cref{fig:coord}). Each
operator acts on the lines of the grid independently, by the ordinary
image $F$:
\begin{equation}\label{eq:columnwise}
 (\Ff P)_j=F(P_j),\quad (\Ff^2P)_j=F^2(P_j),\qquad
 (\Fs Q)_i=F(Q_i),\quad (\Fs^2Q)_i=F^2(Q_i).
\end{equation}
In particular $\abs{\Ff P}=\sum_j\abs{F(P_j)}$ and $\abs{\Fs Q}=\sum_i\abs{F(Q_i)}$.
This is how an inequality about vertex sets, such as strict concavity,
becomes an inequality about sets of pairs: apply it to every column (or
row) and add.

\begin{figure}[ht]
\centering
\begin{tikzpicture}[scale=0.9]
 % relation
 \begin{scope}[xshift=-3.2cm,yshift=-2.2cm]
  \node[vtx] (n1) at (0,1.2) {$1$};
  \node[vtx] (n2) at (-1,0) {$2$};
  \node[vtx] (n3) at (1,0) {$3$};
  \node[vtx] (n4) at (0,-1.2) {$4$};
  \draw[arc] (n1)--(n2); \draw[arc] (n1)--(n3); \draw[arc] (n3)--(n4);
  \node[font=\small] at (0,-2.0) {the relation};
 \end{scope}
 % grid
 \begin{scope}
  \foreach \k in {1,...,4} {
   \node[font=\small] at (\k,0.1) {$\k$};
   \node[font=\small] at (0.2,-\k) {$\k$};
  }
  \node[font=\footnotesize,text=black!60] at (2.5,0.55) {column $j$ (second coordinate)};
  \node[font=\footnotesize,text=black!60,rotate=90] at (-0.35,-2.5) {row $i$ (first coord.)};
  \foreach \i in {1,...,4} \foreach \j in {1,...,4} \node[gridpt] at (\j,-\i) {};
  \node[Lpt,label={[famL!80!black,font=\small]above right:$r$}] (r) at (3,-1) {};
  \node[Rpt,label={[famR!80!black,font=\small]below left:$c$}] (c) at (1,-2) {};
  \node[Zpt,label={[zcol!70!black,font=\small]above right:$z$}] (z) at (3,-2) {};
  \node[circle,draw=famL,thick,inner sep=0pt,minimum size=3.6mm] (r3) at (3,-3) {};
  \node[rectangle,draw=famR,thick,inner sep=0pt,minimum size=3.4mm] (c2) at (2,-2) {};
  \draw[-{Stealth[length=2mm]},famL,thick] (r) -- (z);
  \draw[-{Stealth[length=2mm]},famL,thick] (r) to[bend left=35] (r3);
  \draw[-{Stealth[length=2mm]},famR,thick] (c) -- (c2);
  \draw[-{Stealth[length=2mm]},famR,thick] (c) to[bend right=35] (z);
  \node[font=\small,align=center] at (2.5,-5.0) {$r=(1,3)$: $\Ff\{r\}=\{(2,3),(3,3)\}$\\
   $c=(2,1)$: $\Fs\{c\}=\{(2,2),(2,3)\}$};
 \end{scope}
 % abstract rectangle
 \begin{scope}[xshift=7.8cm,yshift=-0.9cm]
  \node[Hpt,label={[hcol!80!black,font=\small]above left:$h=(p,s)$}] (h) at (0,0) {};
  \node[Lpt,label={[famL!80!black,font=\small]above right:$r=(p,j)$}] (rr) at (2.6,0) {};
  \node[Rpt,label={[famR!80!black,font=\small]below left:$c=(i,s)$}] (cc) at (0,-2.4) {};
  \node[Zpt,label={[zcol!70!black,font=\small]below right:$z=(i,j)$}] (zz) at (2.6,-2.4) {};
  \draw[-{Stealth[length=2mm]},thick,famL] (rr) -- node[right,font=\footnotesize,black] {$p\to i$} (zz);
  \draw[-{Stealth[length=2mm]},thick,famR] (cc) -- node[below,font=\footnotesize,black] {$s\to j$} (zz);
  \draw[black!40,dashed] (h)--(rr); \draw[black!40,dashed] (h)--(cc);
  \node[font=\small,align=center] at (1.3,-4.1) {a conflict and its\\ two ``free'' corners};
 \end{scope}
\end{tikzpicture}
\caption{Left: a relation on $X=\{1,2,3,4\}$. Middle: the grid $X\times X$
(row $=$ first coordinate). The blue point $r$ moves vertically under $\Ff$,
the orange point $c$ moves horizontally under $\Fs$, and their images meet
at the green point $z$. So $r$ and $c$ \emph{conflict}. Right: in general a
conflict between $r=(p,j)$ and $c=(i,s)$ spans a rectangle. The meeting
point $z=(i,j)$ is the opposite corner to the purple point $h=(p,s)$.}
\label{fig:coord}
\end{figure}

\subsection{Conflicts and compatibility}

Throughout, we deal with a \emph{left} family (drawn blue: $P$, $\Rf$, $P'$),
to which $\Ff$ is applied, and a \emph{right} family (drawn orange: $Q$,
$\Cf$, $Q'$), to which $\Fs$ is applied.

\begin{definition}[Conflict, compatibility]\label{def:conflict}
A left member $r=(p,j)$ and a right member $c=(i,s)$ \emph{conflict} if
\[
 p\to i\qquad\text{and}\qquad s\to j.
\]
A left family $P$ and a right family $Q$ are \emph{compatible} if no member
of $P$ conflicts with a member of $Q$.
\end{definition}

\begin{lemma}\label{lem:conflict-geometry}
\begin{enumerate}[label=(\alph*)]
\item $r=(p,j)$ and $c=(i,s)$ conflict if and only if $\Ff\{r\}\cap\Fs\{c\}\ne\varnothing$,
and then $\Ff\{r\}\cap\Fs\{c\}=\{(i,j)\}$.
\item $P$ and $Q$ are compatible if and only if $\Ff P\cap\Fs Q=\varnothing$.
\end{enumerate}
\end{lemma}

\begin{proof}
(a) $\Ff\{r\}$ lies in column $j$ and $\Fs\{c\}$ lies in row $i$, so they can
only meet at $(i,j)$. This point is in $\Ff\{r\}$ if and only if $p\to i$, and
in $\Fs\{c\}$ if and only if $s\to j$. (b) follows from (a), since
$\Ff P=\bigcup_{r\in P}\Ff\{r\}$ and $\Fs Q=\bigcup_{c\in Q}\Fs\{c\}$.
\end{proof}

We say that a left member $r$ \emph{covers} the points of $\Ff\{r\}$, and a
right member $c$ covers the points of $\Fs\{c\}$. So the points covered by a
pair of families $(P,Q)$ are those of $\Ff P\cup\Fs Q$.

\begin{definition}[Conflict points and opposite corners]\label{def:ZH}
For a left family $\Rf$ and a right family $\Cf$ let
\begin{align*}
 \cZ{Z}&=\{(i,j):\ (p,j)\in\Rf\text{ conflicts with }(i,s)\in\Cf\text{ for some }p,s\},\\
 \cH{H}&=\{(p,s):\ (p,j)\in\Rf\text{ conflicts with }(i,s)\in\Cf\text{ for some }i,j\}.
\end{align*}
\end{definition}

So every conflict spans a rectangle with corners $r,z,c,h$ as in
\cref{fig:coord}. The conflict point $z$ is the point both members cover,
and $h$ is the remaining corner. Conversely, $z$ and $h$ determine the
conflict: $r=(h_1,z_2)$ and $c=(z_1,h_2)$.

\subsection{The extremal set-up}

Now let the relation be the arc relation of a nonempty oriented graph $G$
with vertex set $X$, and let $\Delta=\{(v,v):v\in X\}$ be the diagonal.

\begin{definition}[Uncovered points, value, admissible pairs]\label{def:value}
For compatible families $P,Q\subseteq X\times X$ put
\[
 \unc(P,Q)=\abs X^2-\abs{\Ff P}-\abs{\Fs Q},\qquad
 \val(P,Q)=\abs P+\abs Q+\unc(P,Q).
\]
By \cref{lem:conflict-geometry}(b), $\unc(P,Q)$ is the number of grid points
covered by \emph{neither} family. A pair $(P,Q)$ is \emph{admissible} if it
is compatible and $\Delta\subseteq P\cap Q$.
\end{definition}

\begin{lemma}\label{lem:diag-compatible}
$(\Delta,\Delta)$ is admissible. In particular admissible pairs exist, and since
there are finitely many, some admissible pair maximizes $\val$.
\end{lemma}

\begin{proof}
A point $(i,j)\in\Ff\Delta\cap\Fs\Delta$ would require $j\to i$ (from
$(j,j)\in\Delta$) and $i\to j$ (from $(i,i)\in\Delta$). These are opposite
arcs, which an oriented graph does not have.
\end{proof}

\subsection{Why maximizing the value leads to a contradiction}\label{sec:ledger}

Here is the plan of \cref{sec:contradiction}, with the verifications
postponed. Let $(P,Q)$ be an admissible pair of maximum value, write
$\mu=\abs P+\abs Q$ and $d=\unc(P,Q)$, so $\val(P,Q)=\mu+d$.

\begin{enumerate}[label=\textbf{\arabic*.},leftmargin=*]
\item \textbf{Concavity on every line.} Every column $P_j$ and every row $Q_i$ is a
nonempty proper subset of $X$. Applying strict concavity to each and adding
(using~\eqref{eq:columnwise}) gives
\[
 \abs P+\abs{\Ff^2P}<2\abs{\Ff P},\qquad \abs Q+\abs{\Fs^2Q}<2\abs{\Fs Q}.
\]
\item \textbf{Complement.} Let $\Rf=(X\times X)\setminus\Fs^2Q$ and $\Cf=(X\times X)\setminus\Ff^2P$.
Adding the two inequalities shows $\abs\Rf+\abs\Cf>\mu+2d$. This is a gain of
more than $d$ over the old value.
\item \textbf{Repair.} $(\Rf,\Cf)$ contains the diagonal but may have conflicts. The
pruning lemma deletes some members to remove all conflicts. The number
deleted is at most $\abs H$ plus the number of points of $Z$ left
uncovered. We will see that $\abs H\le d$, and the uncovered points number
at most $d'$, the value of $\unc$ for the new pair.
\item \textbf{New value.} The new pair $(P',Q')$ is admissible and
\[
 \val(P',Q')\ \ge\ (\abs\Rf+\abs\Cf)-(d+d')+d'\ >\ \mu+2d-d=\mu+d=\val(P,Q),
\]
contradicting maximality.
\end{enumerate}

\begin{figure}[ht]
\centering
\begin{tikzpicture}[yscale=0.6,xscale=1.0,lab/.style={font=\small,align=center,anchor=north}]
 \def\bw{1.3}
 % levels: mu=4, d=2, d'=1.5 ; R+C = 8.4
 \draw[black!30] (-0.4,0)--(11.1,0);
 \draw[dashed,thick] (-0.4,6)--(11.1,6) node[right,font=\small] {$\mu+d$};
 \fill[famL!45] (0,0) rectangle (\bw,4); \fill[black!18] (0,4) rectangle (\bw,6);
 \node[font=\small] at (0.65,2) {$\mu$}; \node[font=\small] at (0.65,5) {$d$};
 \node[lab] at (0.65,-0.2) {1. old value\\$\val(P,Q)$};
 \fill[famR!45] (2.4,0) rectangle (2.4+\bw,8.4);
 \node[font=\small,align=center] at (3.05,4.2) {$>$\\$\mu{+}2d$};
 \node[lab] at (3.05,-0.2) {2. complement\\$\abs\Rf+\abs\Cf$};
 \fill[red!35] (4.8,4.9) rectangle (4.8+\bw,8.4);
 \node[font=\small,align=center] at (5.45,6.65) {$-\abs H$\\$-q$};
 \node[lab] at (5.45,-0.2) {3. pruning\\cost};
 \fill[zcol!40] (7.2,4.9) rectangle (7.2+\bw,6.4);
 \node[font=\small] at (7.85,5.65) {$+d'$};
 \node[lab] at (7.85,-0.2) {3'. new\\uncovered};
 \fill[famL!25] (9.6,0) rectangle (9.6+\bw,6.4);
 \node[font=\small,align=center] at (10.25,3.2) {new\\value};
 \node[lab] at (10.25,-0.2) {4. $\val(P',Q')$};
 \draw[black!40] (\bw,6)--(2.4,6); \draw[black!40] (2.4+\bw,8.4)--(4.8,8.4);
 \draw[black!40] (4.8+\bw,4.9)--(7.2,4.9); \draw[black!40] (7.2+\bw,6.4)--(9.6,6.4);
\end{tikzpicture}
\caption{The ledger: steps 1--4 above as bar heights. Here $q\le d'$ is the
number of conflict points left uncovered after pruning, so the loss
$\abs H+q\le d+d'$ is more than paid for by the gain and the returned $d'$.
The new bar ends above the dashed line.}
\label{fig:ledger}
\end{figure}

\begin{remark}[Why this particular value]\label{rem:why-value}
If we maximized only $\abs P+\abs Q$, the ledger would give a new size
$>\mu+d-d'$, which need not exceed $\mu$. If we used $\abs P+\abs Q+\lambda\,\unc$
for some $\lambda>0$, the ledger would give $\val_\lambda(P',Q')>\mu+d+(\lambda-1)d'$
against the old value $\mu+\lambda d$. The new value is at least the old one
exactly when $\mu+d+(\lambda-1)d'\ge\mu+\lambda d$, that is, when
$(\lambda-1)(d'-d)\ge0$. Since $d'-d$ can have either sign, this holds for
every $d,d'\ge0$ only when $\lambda=1$. (The strict increase comes from the
strict inequality in the ledger, not from the weight.) So it is essential that the pruning cost is charged partly
to points left uncovered \emph{after} pruning ($d'$), and that this charge
enters with the same weight as the uncovered points in the value.
\end{remark}
